\subsection{A4: the complete order-seven census with nontrivial universal group}\label{sec:A-census}
\begin{proposition}\label{prop:A-small-group-core}
If a finite monoid $M$ has at most eight elements and $U(M)\ne1$, then
$BM\simeq BU(M)$. In particular all its integral homology, in every degree,
is group homology of the explicitly computed finite group $U(M)$.
\end{proposition}
\begin{proof}
Write the minimal ideal in Rees form $I\times H\times J$. The minimal-ideal
argument of Lemma~\ref{A:minimal-group} gives a surjection $H\to U(M)$,
so $|H|\ge2$. If $|I|,|J|\ge2$, the ideal has at least eight elements.
It cannot contain the identity of $M$: a finite simple monoid is a group,
which would give $|I|=|J|=1$. Thus this case requires $|M|\ge9$.
Consequently $|I|=1$ or $|J|=1$. An idempotent $e$ in this ideal then
satisfies $eM=eMe=H$ or $Me=eMe=H$.

Here is the required comparison, including the nonunital inclusion.
Suppose $eM=eMe$; the other case follows by opposite categories.
The homomorphism $r:M\to H$, $r(m)=em$, preserves identities when the
identity of $H$ is $e$, since $eme=em$. Regard $H$ as the full subcategory
on $(\ast,e)$ in the Karoubi envelope $\operatorname{Kar}(M)$.
The inclusion $J:M\to\operatorname{Kar}(M)$ selecting $(\ast,1)$
is an equivalence on classifying spaces by Quillen A. Indeed, for an
idempotent $f$, the comma category $((\ast,f)\mathbin{\downarrow}J)$
has objects $m\in Mf$ and arrows $x:m\to m' $ satisfying $xm=m'$.
The object $f$ and the arrows $m:f\to m$ define a natural transformation
from its constant functor to the identity, so the comma category is contractible. More explicitly, the functors
$M\to\operatorname{Kar}(M)$ selecting $(\ast,1)$ and $r$ followed by the
inclusion of $(\ast,e)$ have the natural transformation with component
$e:(\ast,1)\to(\ast,e)$, because $em=eme$. Also $r|_H=\id_H$.
Thus these maps are inverse after realization. This is precisely the
natural-transformation proof in \cite[Theorem 2.7 and Corollary 2.8]{sweeney2025},
whose complete proof was checked. Hence $BM\simeq BH$; the induced
fundamental-group isomorphism identifies $H$ with $U(M)$.
\end{proof}

\begin{theorem}\label{thm:A-seven-census}
The independent exhaustive search of monoids on seven labelled elements,
with identity fixed to label zero, finds $18\,507\,440$ associative tables.
Exactly $1\,846$ isomorphism-or-opposition classes have nontrivial universal
group. Their labelled multiplicities are as follows. The last columns give,
for every prime dividing the universal group order, the first positive degree
of nonzero homology over that prime field.
\[
\begin{array}{c|r|c|c}
 U(M)&\text{labelled tables}&p&\min\{n>0:H_n(M;\F_p)\ne0\}\\\hline
 C_2&1\,470\,630&2&1\\
 C_3&98\,700&3&1\\
 C_4&14\,040&2&1\\
 C_2\times C_2&4\,920&2&1\\
 C_5&1\,260&5&1\\
 C_6&360&2,3&1,1\\
 S_3&120&2,3&1,3\\
 C_7&120&7&1
\end{array}
\]
Every one of these classes is supplied with its complete multiplication table,
labelled multiplicity, universal-group order and reducing group corner in the
machine-readable output. The $U=1$ labelled count is $16\,917\,290$.
\end{theorem}
\begin{proof}
The program \texttt{c\_enum.cpp} fills each unknown product in turn and
propagates the equations $(ab)c=a(bc)$ whenever three of their four table
entries are known. It rejects a branch only upon an inconsistent forced
entry. Thus every associative table occurs exactly once and no table is
lost. The identity row and column are fixed. The universal group is computed
by identifying every idempotent with the identity and taking the least
multiplicative congruence containing these identifications. Its quotient is
a group: every element of a finite monoid has an idempotent power, so after
these identifications every element has an inverse. The quotient is therefore
exactly universal, not just an approximation to it.

Canonicalization enumerates all permutations within blocks of invariants
which every isomorphism preserves (sizes of principal ideals, fixed points,
idempotence and the index and period of powers), and also the opposite table.
The least resulting string is an exact canonical form. The retained classes
with $|U|=2,3,4,5,6,7$ number $1620,167,52,4,2,1$, respectively.
Their sum is $1846$. The separate homology program rechecks every retained
table's identity and all associativity equations, and reduces successively
at every proper corner satisfying $eMe=eM$ or $eMe=Me$.
It ends in one of the groups in the displayed table, and checks that this
core has exactly the independently computed universal-group order.
Proposition~\ref{prop:A-small-group-core} supplies the homotopy comparison
for every reduction, so these certificates determine all degrees of homology.

For $C_m=\langle g\rangle$ use the complete periodic resolution with maps
$g-1$ and $1+g+\cdots+g^{m-1}$. Tensoring with $\Z$ gives alternating
zero and multiplication by $m$. Hence $H_{2r+1}(C_m;\Z)=\Z/m$ and
$H_{2r}(C_m;\Z)=0$ for $r\ge1$. Tensor products of the two $C_2$
resolutions handle $C_2\times C_2$. These give the degree-one entries.
For $S_3$, the exact integral normalized-nerve computation gives
$H_1=\Z/2$, $H_2=0$, $H_3=\Z/6$; its boundary matrices and Smith
invariants are computed by \texttt{c\_homology.py}. The universal
coefficient theorem then gives first degree one at $p=2$ and degree three
at $p=3$. These are exactly the degrees needed for A4; higher nonvanishing
at both primes also follows from Theorem~\ref{A:bounded}.
\end{proof}
